% !PRINTED 173
\weberpage{173}
If $A$ denotes a substitution of the finite group $S$, we
can derive from $\Theta(x)$ a new function:
\weq{2}{\Theta[A(x)] = c_1' x_1 + c_2' x_2 + \cdots + c_n x_n} % [sic]
whose coefficients $c_i'$ are connected, by \S.\,37, 9., with the original
coefficients $c_i$ through the \gesperrt{transposed substitution corresponding to $A$},
\[ (c') = A_1(c) \]

If now $A, B, C, \dots$ are all the substitutions of the
group $S$, one can form in the same way the functions
\weq{3}{\Theta[A(x)],\ \Theta[B(x)],\ \Theta[C(x)],\ \dots}
which we also denote more briefly by
\weq{4}{\Theta,\ \Theta_1,\ \Theta_2,\ \dots,\ \Theta_{\mu-1}}
where $\mu$ is the degree of the group $S$.

Now one can dispose of the coefficients $c_i$, which up to now
are still entirely arbitrary, in such a way that no two of the functions
(4) become identical with one another (Bd.\ I, \S.\,143, 1.).
But from (3) it follows:

\gesperrt{1. If one applies simultaneously in the $\mu$ functions
(4) any substitution from $S$, the
totality of these functions does not change,
but they suffer only a permutation.}

From this it follows further:

\gesperrt{2. Every symmetric function of the quantities (4) is
an absolute invariant of the group $S$.}

We remark further that in place of the function $\Theta$
one could put any other function, even a non-linear, indeed a fractional
or inhomogeneous one, provided only that the $\mu$ functions
(4) are different from one another.

If we now form the product:
\weq{5}{\begin{aligned}
\Phi(t) &= (t-\Theta)(t-\Theta_1)\dots(t-\Theta_{\mu-1}) \\
&= t^\mu + A_1 t^{\mu-1} + A_2 t^{\mu-2} + \cdots + A_\mu,
\end{aligned}}
which is an integral rational function of degree $\mu$ in $t$, then
the coefficients $A_1, A_2, \dots, A_\mu$ are, by 2., invariants of the
group $S$, and the quantities $\Theta, \Theta_1, \dots, \Theta_{\mu-1}$ are the roots
of the equation
\weq{6}{\Phi(t) = 0.}
